% Reconstructed from the retained accepted proof after the execution-environment
% loss of 16 November 2026. This is not a byte-identical recovery of the old file.
\subsection*{The common direction survives every rational unipotent level}

\statement{B.66}{Proved}{A nonzero common prounipotent line survives all rigid tests}
Fix a prime $\ell$.  Let $\mathfrak p_\ell$ be the Lie
algebra of the $\mathbb Q_\ell$-prounipotent completion
of the actual based group $P=P\Gamma_{0,12}$, and let
$\mathfrak h_{T,\ell},\mathfrak h_{5,\ell}$ be those
of the two actual source groups $H_T,H_5$.
There are compatible elements in their commutator
filtrations
\[
 A_\ell\in F^2\mathfrak h_{T,\ell},\qquad
 B_\ell\in F^2\mathfrak h_{5,\ell}
\]
whose images in $\mathfrak p_\ell$ are the same element
$W_\ell$, such that
\[
 \operatorname{gr}_2W_\ell=w\ne0,\qquad
 a(W_\ell)=W_\ell\quad(a\in A_5),\qquad
 f_S(W_\ell)=0\quad(S\in\mathscr S_7).
                                                               \tag{B.66.1}
\]
Here $w$ is B.59's primitive integral vector, viewed over
$\mathbb Q_\ell$, and every $f_S$ is the full
prounipotent forgetful homomorphism.
Consequently the nonzero one-parameter subgroup
\[
                \exp(\mathbb Q_\ell W_\ell)
                   \cong\mathbb G_a(\mathbb Q_\ell)
                                                               \tag{B.66.2}
\]
is contained in both actual source prounipotent images,
is fixed by $A_5$, and is killed by all rigid-seven tests
and their smaller rigid factors.
This is a result about rational prounipotent completions.
It does not construct a nonidentity element in a compact
pro-$\ell$ completion or in the full residual group
$\mathcal K$.

\begin{proof}
\emph{A strictly pointed arithmetic diagram.}
Use B.61's actual ordered common base configuration,
not two separately chosen abstract identifications.
With $\varphi=(1+\sqrt5)/2$, its source parameters are
\[
 u_0=i\bigl(\varphi+\sqrt{\varphi+2}\bigr),\qquad
 v_0=-\varphi^2.
                                                               \tag{B.66.3}
\]
The tetrahedral and cyclic labelled maps agree at these
points after normalizing labels $2,4,6$ to
$\infty,0,1$.  B.61 verifies those equalities in
$\mathbb Q(i,\sqrt5,\sqrt{\varphi+2})$ and supplies the
common topological transports.
Enlarge this field to a finite number field $K$ over
which every puncture of both rational source curves and
every coefficient of the two maps is defined.
We obtain based $K$-morphisms
\[
 (S_T,u_0)\longrightarrow(M_{0,12},b)
        \longleftarrow(S_5,v_0),
                                                               \tag{B.66.4}
\]
together with all the based forgetting maps.
The ordered space $M_{0,12}$ is a fine moduli scheme.
Every relabelling in the actual icosahedral $A_5$ fixes
the point $b$ literally, because its permutation is
realized by the corresponding automorphism of that
marked sphere.  These $K$-automorphisms fixing $b$
give the actual based $A_5$ action used in B.59.

Choose compatible geometric points over the displayed
$K$-points.  The geometric--arithmetic exact sequence
has a section supplied by each rational point.
The functorial homomorphisms of fundamental groups
commute with these sections exactly.  In particular
the action of $G_K$ on the geometric groups commutes
with every map in (B.66.4), every forgetting map, and
the actual $A_5$ action.  There is no choice of an
uncontrolled outer representative in this assertion.
The required exact functoriality and the exact sequence
are respectively Theorem~58.6.2(3) and Lemma~58.14.3
of the Stacks Project \citep{DE_Stacks2026}.
The comparison with the based topological groups and
their $\mathbb Q_\ell$-unipotent completions is natural;
the discrete, pro-$\ell$ and full profinite versions
give the same unipotent completion
\citep[Appendix A, Theorem A.6, Lemma A.7 and
Corollary A.11]{DE_HainMatsumoto2003}.

\emph{Scalar action on every graded piece.}
For a rational curve whose punctures are all $K$-rational,
the first homology is a direct sum of copies of
$\mathbb Q_\ell(1)$, as is seen from the Kummer maps
$t-a$.  For ordered $M_{0,n}$, normalization of three
labels gives the units $z_i$, $z_i-1$ and $z_i-z_j$;
their Kummer classes and the usual pair-linking
description give the same assertion for first homology.
Thus $\sigma\in G_K$ acts on each first homology here
as the scalar $\chi_\ell(\sigma)$.
The lower-central associated graded Lie algebra is
generated by degree one.  Therefore its degree-$d$
piece has scalar action $\chi_\ell(\sigma)^d$.

Since $K$ is a number field, the cyclotomic image is
open in $\mathbb Z_\ell^\times$.
Choose $\sigma\in G_K$ for which
$\alpha=\chi_\ell(\sigma)$ has infinite order, and
write $\Phi$ for its action on any of the Lie algebras.
On every finite class-$c$ quotient, the factors below
commute and successively raise the filtration, so
\[
              \prod_{d=1}^c(\Phi-\alpha^d)=0.
                                                               \tag{B.66.5}
\]
The roots $\alpha,\ldots,\alpha^c$ are distinct.
Consequently $\Phi$ is semisimple on this quotient,
with its degree-two projector explicitly given by
\[
 E_{2,c}(\Phi)=
       \prod_{\substack{1\le d\le c\\d\ne2}}
            \frac{\Phi-\alpha^d}{\alpha^2-\alpha^d}
                 \qquad(c\ge2).
                                                               \tag{B.66.6}
\]
The $\alpha^2$ eigenspace maps isomorphically to
$\operatorname{gr}_2$: an eigenvector has zero image
in degree one and cannot have its first nonzero term
in any degree greater than two.  Conversely (B.66.6)
applied to a lift of a degree-two class supplies its
unique eigenlift.  The construction is functorial for
all the strictly equivariant maps above and is compatible
with passage from class $c+1$ to class $c$.

\emph{Compatible preimages in the actual sources.}
Let $a\in H_T'$ and $b_0\in H_5'$ be B.60's actual
finite words.  Their ambient degree-two images are
$2w$ and $w$, respectively.  In each source class-$c$
Lie algebra set
\[
 A_{\ell,c}=\tfrac12 E_{2,c}(\Phi_T)\log a,
 \qquad
 B_{\ell,c}=E_{2,c}(\Phi_5)\log b_0.
                                                               \tag{B.66.7}
\]
The factor $1/2$ is legitimate also when $\ell=2$,
since the coefficient field is $\mathbb Q_\ell$.
Their ambient images have degree-two term $w$ and
are both $\alpha^2$ eigenvectors, so uniqueness of
the eigenlift gives a common image $W_{\ell,c}$.
The source elements in (B.66.7) are themselves
compatible in $c$.  Taking their inverse limits
therefore supplies actual elements $A_\ell,B_\ell$
in the two complete source Lie algebras and a common
ambient image $W_\ell$.  This establishes membership
in the two full source images, rather than merely
membership in both images at every finite stage.

By B.59, $A_5$ fixes $w$ and each rigid-seven map
kills $w$ in degree two.  Since the $A_5$ action
commutes with $\Phi$, the vector $a(W_{\ell,c})$
is another eigenlift of $w$ and hence equals
$W_{\ell,c}$.  A rigid image of $W_{\ell,c}$ is an
$\alpha^2$ eigenvector in filtration at least three,
so it is zero.  These identities are compatible
in $c$, and prove (B.66.1).
Exponentiation, which is defined on every rational
nilpotent quotient and is compatible with inverse
limits, proves (B.66.2).  Its injectivity follows
already from the nonzero degree-two vector $w$.

\emph{The integral limitation.}
The proof uses divisions by
$\alpha^2-\alpha^d$.  Taking $\alpha$ in a sufficiently
small principal subgroup if necessary, their
$\ell$-adic valuations are unbounded along
$d=2+\ell^j$.
No uniform denominator bound for the compatible
eigenlifts has been established.  Clearing denominators
separately at each finite nilpotent class need not give
one nonzero multiplier that works at every class.
Possible cancellations are not ruled out, but they
are not proved either.  Thus the rational line does
not supply a compact pro-$\ell$ element or an element
of $\mathcal K$.  It proves that all rational
prounipotent versions of these existing common-image
and rigid-test constraints retain the nonzero common
direction.
\end{proof}
